% 11_proofs.tex -- appendix
\section{Proofs}
\label{app:proofs}

\subsection{Proof of Theorem~\ref{thm:detection}}

Let
\[
  Z_i = \mathbf{T}^{\mathrm{clk}}_{\sigma_k(t_i),\,\sigma_k(t_{i+1})}
\]
be the validity indicator at step $i$.

\paragraph{Null mean.}
Under the random oracle, $\sigma_k$ is uniform over $\mathcal{S}$ and
independent across distinct tokens. Hence $Z_i \sim
\mathrm{Bernoulli}(1/S)$ for random input, and
\[
  \mathbb{E}[\phi \mid \mathrm{null}] = \tfrac{1}{S}.
\]

\paragraph{Watermark mean.}
At a gated position, Algorithm~\ref{alg:embed} sets
$\sigma_k(t_{i+1}) = (\sigma_k(t_i) + 1) \bmod S$, so $Z_i = 1$
deterministically. At an ungated position, $Z_i \sim
\mathrm{Bernoulli}(1/S)$. Averaging,
\[
  \mathbb{E}[\phi \mid \mathrm{WM}]
    = \rho \cdot 1 + (1-\rho) \cdot \tfrac{1}{S}
    = \tfrac{1}{S} + \rho \tfrac{S-1}{S}.
\]

\paragraph{Null variance.}
The null variance is $\mathrm{Var}[\phi \mid \mathrm{null}] =
(1/S)(1-1/S)/(n-1)$. This follows from the pairwise zero-covariance
argument in the proof of Proposition~\ref{prop:kregular} below,
specialised to $k = 1$: clockwork has $|N^-(s')| = 1$, so
column-regularity holds trivially.

\paragraph{Detector $z$-score.}
The signed mean gap is $\rho(S-1)/S$. Dividing by the null SD,
\begin{align*}
  z(\rho, S, n)
    &= \frac{\rho(S-1)/S}{\sqrt{(1/S)(1-1/S)/(n-1)}} \\
    &= \rho\sqrt{(S-1)(n-1)}.
\end{align*}
This is the full mean-gap $z$ reported by Algorithm~\ref{alg:detect}.
The null $z$ \emph{at the midpoint threshold}
$\tau_{\mathrm{mid}} = 1/S + \tfrac{\rho}{2}(S-1)/S$ is half this
quantity, $(\rho/2)\sqrt{(S-1)(n-1)}$, since
$\tau_{\mathrm{mid}} - 1/S$ is half the mean gap.

\paragraph{FPR bound.}
Adjacent indicators $Z_i, Z_{i+1}$ share token $t_{i+1}$, so
$\{Z_i\}$ are not jointly i.i.d.\ and Hoeffding does not apply
directly. Split into
\[
  Z_1, Z_3, Z_5, \ldots
  \qquad\text{and}\qquad
  Z_2, Z_4, Z_6, \ldots,
\]
two disjoint subsequences of $m = \lfloor (n-1)/2 \rfloor$ terms,
each i.i.d.\ $\mathrm{Bernoulli}(1/S)$. Hoeffding on each plus union
bound gives, for $\tau > 1/S$,
\[
  \Pr(\phi \geq \tau)
    \;\leq\; 2 \exp\!\bigl(-2 m\, (\tau - 1/S)^2\bigr).
\]

\paragraph{Calibration.}
The Hoeffding bound is conservative. Pairwise zero covariance plus
boundedness of $Z_i$ in $[0,1]$ yield a CLT-style Gaussian limit for
$\phi$, so we adopt the Gaussian convention.
At the midpoint threshold, the null $z$-score at
$\tau_{\mathrm{mid}}$ is $(\rho/2)\sqrt{(S-1)(n-1)}$. Requiring this
to exceed $z_\alpha$ for $\mathrm{FPR} \leq \alpha$ gives
\[
  \tfrac{\rho}{2}\sqrt{(S-1)(n-1)} \;\geq\; z_\alpha,
\]
i.e.,
\[
  S \;\geq\; \frac{4 z_\alpha^2}{\rho^2(n-1)} + 1.
\]
Taking ceiling yields~\eqref{eq:calibration}. \hfill$\square$

\subsection{Proof of Proposition~\ref{prop:kregular}}

Let $T \in \{0,1\}^{S\times S}$ be $k$-regular, and write
\[
  N^+(s) := \{s' : T(s,s') = 1\},
  \qquad
  N^-(s') := \{s : T(s,s') = 1\}
\]
for the out- and in-neighbourhoods. By $k$-regularity,
$|N^+(s)| = |N^-(s')| = k$ for every $s, s'$.

\paragraph{Null mean \eqref{eq:kreg-null}.}
Under the random oracle, $\sigma_k(t_i)$ and $\sigma_k(t_{i+1})$
are i.i.d.\ uniform on $\mathcal{S}$. Conditioning on
$\sigma_k(t_i)$ and applying row-regularity,
\[
  \Pr[X_i = 1]
    \;=\; \mathbb{E}_{\sigma_k(t_i)}\!\left[
      \tfrac{|N^+(\sigma_k(t_i))|}{S}
    \right]
    \;=\; \tfrac{k}{S}
    \;=\; p_0.
\]

\paragraph{Watermark mean \eqref{eq:kreg-wm}.}
At a gated position, Algorithm~\ref{alg:embed} restricts the argmax
to $\bigcup_{s' \in N^+(s)} \mathcal{V}_{s'}$, so
$\sigma_k(t_{i+1}) \in N^+(\sigma_k(t_i))$ and $X_i = 1$. At an
ungated position, $X_i \sim \mathrm{Bernoulli}(p_0)$. Averaging,
\[
  \mathbb{E}[\phi \mid \mathrm{WM}]
    \;=\; \rho + (1-\rho) p_0
    \;=\; p_0 + \rho(1 - p_0).
\]

\paragraph{Null variance \eqref{eq:kreg-var}.}
Write $\phi = (n-1)^{-1} \sum_i X_i$. Non-adjacent pairs
$(X_i, X_j)$ are independent: they share no tokens. For adjacent
pairs $(X_i, X_{i+1})$ which share $\sigma_k(t_{i+1}) = s'$,
\begin{align*}
  \Pr[X_i = 1 \mid \sigma_k(t_{i+1}) = s']
    &= \tfrac{|N^-(s')|}{S} = p_0, \\
  \Pr[X_{i+1} = 1 \mid \sigma_k(t_{i+1}) = s']
    &= \tfrac{|N^+(s')|}{S} = p_0.
\end{align*}
\emph{Column-regularity is essential}: without it, the first
probability would depend on $s'$.

The conditional independence
$\sigma_k(t_i) \perp \sigma_k(t_{i+2}) \mid \sigma_k(t_{i+1})$
holds when $t_i \neq t_{i+2}$, since $\mathcal{H}(k \,\Vert\, \cdot)$ is
queried on distinct inputs. Coincidences $t_i = t_{i+2}$ contribute
$O(1/|\mathcal{V}|)$ to the pair covariance under an i.i.d.\ uniform
null on tokens, negligible for realistic vocabularies.
Then
\[
  \Pr[X_i = 1, X_{i+1} = 1]
    = \mathbb{E}_{s'}[p_0^2]
    = p_0^2,
\]
so $\mathrm{Cov}(X_i, X_{i+1}) = 0$. All pairwise covariances
vanish, giving
\[
  \mathrm{Var}(\phi)
    = \frac{p_0(1-p_0)}{n-1}.
\]

\paragraph{$z$-score \eqref{eq:kreg-z}.}
Dividing the signed mean gap $\rho(1-p_0)$ by the null SD,
\begin{align*}
  z &= \frac{\rho(1-p_0)}{\sqrt{p_0(1-p_0)/(n-1)}} \\
    &= \rho \sqrt{\frac{(1-p_0)(n-1)}{p_0}}.
\end{align*}
For $k = 1$ (clockwork) $p_0 = 1/S$ and this reduces to
$\rho\sqrt{(S-1)(n-1)}$, the full mean-gap $z$ of
Theorem~\ref{thm:detection}. Halving the gap (midpoint threshold)
recovers the calibration form
$(\rho/2)\sqrt{(S-1)(n-1)}$ used in~\eqref{eq:calibration}.

\paragraph{Post-attack mean and critical rate
\eqref{eq:kreg-atk}--\eqref{eq:kreg-delta}.}
Let $M_i \sim \mathrm{Bernoulli}(\delta)$ indicate modification of
$t_i$. Modified tokens receive fresh uniform $\sigma_k$ by the
random oracle.

For each pair $(i, i+1)$, condition on $(M_i, M_{i+1})$:
\begin{itemize}[nosep, leftmargin=1.2em]
  \item $(0,0)$: contributes the unattacked mean
    $p_0 + \rho(1-p_0)$.
  \item $(0,1)$: $t_{i+1}$ fresh; column-regularity gives
    $\mathbb{E}[X_i] = p_0$.
  \item $(1,0)$: symmetric; row-regularity gives $\mathbb{E}[X_i]
    = p_0$.
  \item $(1,1)$: both fresh; $\mathbb{E}[X_i] = p_0$.
\end{itemize}
Combining with weights $(1-\delta)^2$ and $1-(1-\delta)^2$:
\begin{align*}
  \mathbb{E}[\phi \mid \mathrm{atk}]
    &= (1-\delta)^2\bigl(p_0 + \rho(1-p_0)\bigr)
       + \bigl(1 - (1-\delta)^2\bigr) p_0 \\
    &= p_0 + \rho (1-\delta)^2 (1-p_0).
\end{align*}

At the midpoint threshold
$\tau_{\mathrm{mid}} = p_0 + \tfrac{\rho}{2}(1-p_0)$, detection
is defeated iff
\[
  \rho(1-\delta)^2(1-p_0) \;<\; \tfrac{\rho}{2}(1-p_0),
\]
i.e., $(1-\delta)^2 < 1/2$. Solving,
\[
  \delta^\star_{\mathrm{mid}}
    \;=\; 1 - 1/\sqrt{2},
\]
independent of $S, \rho, n, k$. \hfill$\square$

\subsection{Proof of Theorem~\ref{thm:quality}}

\paragraph{KL formula.}
At a gated position with current state $s$, the renormalised MCL
distribution is
\[
  P_{\mathrm{MCL}}(t)
    = \begin{cases}
        p(t)/Z_s, & t \in \mathcal{V}_{(s+1) \bmod S}, \\
        0, & \text{otherwise},
      \end{cases}
\]
with $Z_s = p(\mathcal{V}_{(s+1) \bmod S})$. Direct computation:
\begin{align*}
  D_{\mathrm{KL}}(P_{\mathrm{MCL}} \,\|\, p)
    &= \sum_{t \in \mathcal{V}_{(s+1) \bmod S}}
        \frac{p(t)}{Z_s} \log \frac{1}{Z_s} \\
    &= \log(1/Z_s).
\end{align*}

\paragraph{Jensen lower bound.}
Under hash uniformity, $\Pr[t \in \mathcal{V}_{(s+1) \bmod S}] = 1/S$
for every fixed $t$, so by linearity
\[
  \mathbb{E}[Z_s]
    = \sum_{t \in \mathcal{V}} p(t) \cdot \tfrac{1}{S}
    = \tfrac{1}{S}.
\]
The map $z \mapsto \log z$ is concave; Jensen gives
$\mathbb{E}[\log Z_s] \leq \log \mathbb{E}[Z_s] = -\log S$, hence
\[
  \mathbb{E}[D_{\mathrm{KL}}]
    \;=\; \mathbb{E}[\log(1/Z_s)]
    \;\geq\; \log S,
\]
with equality iff $Z_s$ is constant in the key. Averaging over
positions with gate rate $\rho$,
\[
  \mathbb{E}[D_{\mathrm{KL}}]_{\mathrm{tok}}
    \;\geq\; \rho \log S.
\]
The bound is tight at high-entropy positions (where $p$ is close to
uniform and $Z_s \approx 1/S$), and loose at low-entropy positions
(where $Z_s$ varies sharply with the choice of partition). \hfill$\square$

\subsection{Proof of Theorem~\ref{thm:robustness}}

Let $M_i \sim \mathrm{Bernoulli}(\delta)$ indicate modification of
$t_i$, i.i.d.\ across $i$. We treat the fully-gated case $\rho = 1$
first.

\paragraph{Decomposition for $\rho = 1$.}
Pre-attack, $X_i = 1$ deterministically. Post-attack, decompose on
$(M_i, M_{i+1})$:
\begin{itemize}[nosep, leftmargin=1.2em]
  \item $(0,0)$: valid w.p.\ $1$. Weight $(1-\delta)^2$.
  \item $(0,1)$: $\sigma_k(t_{i+1})$ fresh-uniform; valid w.p.\
    $1/S$. Weight $(1-\delta)\delta$.
  \item $(1,0)$: symmetric. Weight $\delta(1-\delta)$, prob.\ $1/S$.
  \item $(1,1)$: both fresh-uniform. Weight $\delta^2$, prob.\ $1/S$.
\end{itemize}
Summing,
\[
  \mathbb{E}[\phi^{\mathrm{atk}} \mid \rho = 1]
    \;=\; (1-\delta)^2 + \tfrac{\delta(2-\delta)}{S}.
\]

\paragraph{General $\rho$.}
For $\rho < 1$, the $(0,0)$ pre-attack indicator is
$g_i + (1-g_i)/S$, with $\mathbb{E}_g[\cdot] = \rho + (1-\rho)/S$.
The non-$(0,0)$ contributions remain $1/S$: the fresh-token state is
uniform regardless of the gate indicator. Combining,
\begin{align*}
  \mathbb{E}[\phi^{\mathrm{atk}}]
    &= (1-\delta)^2 [\rho + (1-\rho)/S]
       + (1-(1-\delta)^2)/S \\
    &= \tfrac{1}{S} + \rho(1-\delta)^2 \tfrac{S-1}{S}.
\end{align*}
This matches the $k=1$ specialisation of
Proposition~\ref{prop:kregular}'s post-attack identity.

\paragraph{Critical rate.}
At the midpoint threshold
$\tau = 1/S + \tfrac{\rho}{2}(S-1)/S$, the attack succeeds when
\[
  \rho(1-\delta)^2 \tfrac{S-1}{S}
    \;<\;
  \tfrac{\rho}{2} \tfrac{S-1}{S},
\]
i.e., $(1-\delta)^2 < 1/2$. Solving,
$\delta^\star = 1 - 1/\sqrt{2}$, independent of $S, \rho, n$. \hfill$\square$

\subsection{Proof of Theorem~\ref{thm:security}}

The state assignment $\sigma_k(t) = \mathcal{H}(k \,\Vert\, t) \bmod S$
uses $\mathcal{H}$ modelled as a random oracle. Without~$k$, any polynomial-time
adversary $\mathcal{A}$ can evaluate $\sigma_k$ only by querying
the oracle on $k' \,\Vert\, t$ for $k' \neq k$, which yields fresh
independent outputs.

\paragraph{Key recovery.}
Assume $k$ has min-entropy $\lambda$. Since the oracle's output
on $k \,\Vert\, t$ is independent of $\mathcal{A}$'s view, a query
budget $q$ yields recovery probability at most
\[
  q \cdot 2^{-\lambda}.
\]

\paragraph{State prediction.}
For any token $t^\star$ that $\mathcal{A}$ has not previously
queried with the correct key, $\sigma_k(t^\star)$ is uniform on
$\mathcal{S}$ in $\mathcal{A}$'s view. The per-prediction success
probability is at most $1/S + \mathrm{negl}(\lambda)$, and the
\emph{advantage} over the uniform $1/S$ baseline is
$\mathrm{negl}(\lambda)$.

\paragraph{Statistical detection (restricted threat model).}
Consider an adversary $\mathcal{A}$ without oracle access, without
LM access, and without reference text. The null hypothesis
$\mathbf{t}^r$ is i.i.d.\ uniform over $\mathcal{V}$ (as in
Theorem~\ref{thm:detection}). Under this null, any statistic
measurable only in token IDs that achieves non-trivial advantage
must extract $\sigma_k$-dependent structure from $\mathbf{t}^w$.
Without oracle access, $\mathcal{A}$ has at most
$\mathrm{negl}(\lambda)$ advantage on each $\sigma_k$ prediction,
so the distinguishing advantage inherits the same bound.

An adversary \emph{with} LM access can compute the perplexity gap
implied by Theorem~\ref{thm:quality} and is not covered by this
theorem. MCL is a distorting watermark, not an undetectable one.
\hfill$\square$

\subsection{Proof of Theorem~\ref{thm:hdd}}

Under the null, $X_i$ is $\mathrm{Bernoulli}(1/S)$ marginally with
pairwise zero covariance (Proposition~\ref{prop:kregular}). Under
the alternative, $X_i$ is $\mathrm{Bernoulli}(\pi_i)$ with
\[
  \pi_i = g_i + (1 - g_i)/S.
\]

\paragraph{1-dependence and LAN.}
Under the random oracle on distinct tokens, $\{X_i\}$ is a
1-dependent sequence: $X_i$ and $X_{i+1}$ share token $t_{i+1}$,
while $(X_i, X_j)$ for $|i-j| \geq 2$ are independent. Pairwise
covariance vanishes (Prop.~\ref{prop:kregular}), so the central
limit theorem for 1-dependent sequences gives the standardised
product-form LLR the same Gaussian limit as the true LLR. The test
below is therefore \emph{asymptotically} Neyman--Pearson optimal
in the Le~Cam LAN sense. We adopt the convention $0 \log 0 = 0$
for gated $X_i$.

\paragraph{Likelihood ratio.}
Under product hypotheses,
\begin{align*}
  \Lambda
    &= \sum_i \log
        \frac{\pi_i^{X_i}(1-\pi_i)^{1-X_i}}
             {(1/S)^{X_i}(1-1/S)^{1-X_i}} \\
    &= \sum_i X_i \log(\pi_i S)
       + (1-X_i) \log\tfrac{1-\pi_i}{1-1/S},
\end{align*}
matching the theorem statement. Asymptotic Neyman--Pearson
optimality follows from the Neyman--Pearson lemma applied to the
product hypothesis combined with the 1-dependent CLT above.
\hfill$\square$

\subsection{Proof of Theorem~\ref{thm:healing}}

Let $\mathcal{T}_k \subseteq \mathcal{V}$ denote the tokens the
adversary has queried the oracle on with the correct key; the
image $\sigma_k(\mathcal{T}_k)$ is its known-states set.

\paragraph{Oracle-blind contribution.}
For any token $t' \notin \mathcal{T}_k$ that the adversary
introduces, $\sigma_k(t')$ is uniform on $\mathcal{S}$ by the
random-oracle property. Each modified position contributes an
indicator $Z_j \sim \mathrm{Bernoulli}(1/S)$ in the validity count,
independently of the adversary's strategy.

\paragraph{Modified-pair mass.}
The fraction of pairs touching at least one modification is
\[
  1 - (1-\delta)^2 \;=\; \delta(2-\delta).
\]
Each contributes at least $1/S$ in expectation, so the expected
mass contributed by modified pairs satisfies
\[
  \text{(modified-pair mass)}
    \;\geq\; \Sigma(\delta)
    \;:=\; \tfrac{1}{S}\delta(2-\delta),
\]
with equality when every introduced token is unqueried.

\paragraph{Strategy independence.}
An adversary with oracle access on $\mathcal{T}_k$ can craft pairs
of queried tokens satisfying $T$ deterministically, raising
$\phi$ above $\Sigma(\delta)$. This helps detection rather than
hurting it. The lower-bound direction $\geq 1/S$ per modified
pair is therefore strategy-free. \hfill$\square$
